\SourceLeaf{244}%
1829---1830 is hardly likely to be found in Germany. I can, however, give you the title of another work, which you can no doubt obtain very easily through a bookseller and which is only 4 printing sheets long.

\begin{center}
»Geometrische Untersuchungen zur Theorie der Parallel-Linien\\
von \textsc{Nicolaus} \textsc{Lobatschewsky} Kais. Russischem Staatsrath \emph{etc.}\\
Berlin 1840 in der G. Finckeschen Buchhandlung.«
\end{center}

I remember reading a very dismissive review of this book in \textsc{Gersdorf}'s Repertorium at the time. To any reader with some knowledge of the subject, the review bore the stamp of someone wholly ignorant. Since I have had the chance to examine this short work myself, I must judge it \textls{very favorably}. In particular, it has much more elegance and precision than \textsc{Lobachevsky}'s longer papers. Those papers are more like a tangled forest through which it is hard to find a way and gain an overview without first becoming acquainted with every tree individually.

On the experimental bound cited in \textsc{Crelle} 17 p. 303, I have found nothing in the work of 1840, and I shall therefore probably have to make up my mind to write about it to Mr \textsc{Lobachevsky} himself. I arranged for his admission as a corresponding member of our Society a year ago. Perhaps he will then send me the Kazan Messenger. It is possible, though, that something about it appears in the subsequent 1836---1838 volumes of the Kazan Learned Writings, which also contain long papers by \textsc{Lobachevsky}. I possess these volumes, but for the reason mentioned in my previous letter I have not yet become better acquainted with them.

Incidentally, in his letter thanking me for his admission to the Society, \textsc{Lobachevsky} wrote at the time that his many administrative duties (he seems to be Rector perpetuus of the university) had now drawn him entirely away from scientific work.

I close for today with the request that you soon gladden me with a few lines.

\begin{flushright}
Yours faithfully and devotedly\\
\textsc{C. F. Gauss}.
\end{flushright}

\par\bigskip
